\paragraph{一处变体。}\label{sub:a-variation}
我们现在对前面的构造稍作修改，把\trStrip{} $Z$ 换成\trUpperHalfPlane{} $H$。把
$\scr{P}^-\,Gr(V)$ 记为由满足下列两条性质的路径
$\lambda\colon[0;1]\to Gr(V)$ 组成的空间：
\begin{itemize}
\item $\lambda(0)$ 与 $\lambda(1)$ \trTransverse{}相交。
\item 数对 $(\lambda,\lambda(1))$——其中第二个分量指取值为 $\lambda(1)$ 的
常值路径——在 $s=1$ 处具有\trNegativeDefinite{}的\trCrossingForm。
\end{itemize}
该空间上带着一个自然的\trLocallyConstant{}函数 $I_H$ 和一个\trLineBundle{}
$\delta_H$。为了把它化归到先前的设置，注意由上面的第二个条件，存在 $s'<1$ 使得
$\lambda(s)\pitchfork\lambda(1)$ 对一切 $s\in[s';1)$ 成立；特别地，把
$(\lambda,\lambda(1))$ 限制到 $[0;s']$ 上便给出 $\scr{P}^2\,Gr(V)$ 的一个元素。
可以选取 $s'$ 使得它\trSmooth{}地依赖于 $\lambda$，此时限制便定义了一个映射
$\scr{P}^-\,Gr(V)\to\scr{P}^2\,Gr(V)$。取 $I_H$、$\delta_H$ 为 $I_Z$、$\delta_Z$
沿这一映射的\trPullback。更明确地说，以 $q_\lambda$ 代替 $q_{\lambda,\lambda(1)}$
来写，对\trGeneric{}的路径 $\lambda\in\scr{P}^-\,Gr(V)$ 我们有
\begin{equation}\label{eq:halfplane-index-formula}
\begin{gathered}
I_H(\lambda)=\sum_{s<1}\mathrm{sign}(q_\lambda(s)),\\
(\delta_H)_\lambda=\bigotimes_{s<1}
\bigl(\lambda(s)\cap\lambda(1)\bigr)^{\mathrm{sign}(q_\lambda(s))}
\end{gathered}
\end{equation}
在\trIndex{}理论这一侧，如前选取一个函数 $\psi$，并考虑 $E=H\times V$ 上具%
\trBoundaryCondition{} $F_\lambda\subset\partial H\times V=\RR\times V$、
$(F_\lambda)_s=\lambda_{\psi(s)}$ 的 $\bar\partial$-算子。由于 $\lambda$ 两端点的%
\trTransversality，它是\trNondegenerate{}的。将它记作 $D_{H,\lambda}$，则有：
\begin{lemma}\label{thm:halfplane-index-determinant}
$I_H(\lambda)=\mathrm{index}(D_{H,\lambda})$；并且 $(\delta_H)_\lambda$
\trCanonical\trIsomorphic{}于\trDeterminantLine{} $\det(D_{H,\lambda})$。
\end{lemma}
\begin{proof}
我们从一个简单的例子开始，即形如
\begin{equation}\label{eq:model-path-loop}
\lambda(s)=\exp(\alpha s I_V)\lambda(0)
\end{equation}
的路径，其中 $\lambda(0)$ 是某个\trLagrange{}子空间，$\alpha\in(-\pi;0)$ 是一个角。
此时没有任何交点进入\eqref{eq:halfplane-index-formula}，因此引理的陈述是说
$D_{H,\lambda}$ 的\trIndex{}为零，并且自带一个取定的\trIsomorphism
$\det(D_{H,\lambda})\cong\RR$。首先，\trIndex{}不依赖于 $\alpha$ 的选取，故不妨设
$\alpha=-\pi/2$。在这种情形下，可以把 $D_{H,\lambda}$ 的拷贝\trGluing{}成%
\trDisk{}上的一个算子 $D_{D,\rho}$，其关联的\trLagrange{}子空间\trLoop{}满足
$\mu(\rho)=-n$，于是\eqref{eq:delbar-gluing} 与
\cref{thm:det-line-bundle-loop-space} 合起来给出
$2\,\mathrm{index}(D_{H,\lambda})=\mathrm{index}(D_{D,\rho})=0$。
第二部分是要证明 $D_{H,\lambda}$ 为\trInjective{}的。取\trDirectSummand{}可以把
问题化到最小的秩 $n=1$，此时\cref{thm:winding-injectivity}适用；在该引理的记号下，
我们的路径 $\lambda$ 可以补全成一个满足 $\mu(\rho)=0$ 的\trLoop{} $\rho$，而我们的%
\trSurface{} $H$ 满足 $|\Sigma^-|=1$（显然，由于这是一个特别简单的情形，也不难想出
其他同样可行的办法）。
至于一般情形，$\scr{P}^-\,Gr(V)$ 中的任何路径都可以在不改变 $I_Z$ 或 $\delta_Z$
的前提下\trDeformation{}成在某条 $[s';1]$（$s'<1$）上满足
\eqref{eq:model-path-loop} 的路径，故不妨只考虑已具此性质者。把
$(\lambda,\lambda(1))$ 限制到 $[0;s']$ 上给出 $\scr{P}^2\,Gr(V)$ 的一个元素，从而
给出 $Z$ 上的一个 $\bar\partial$-算子，其\trIndex{}与\trDeterminantLine{}由
\cref{thm:path-space-det-line} 给出，因而与\eqref{eq:halfplane-index-formula}
相合。另一方面，$H$ 上关联于 $\lambda|_{[s';1]}$ 的算子则落入先前的论证。把 $H$
的端连接到 $Z$ 的出端，便可将这两块\trGluing{}起来，其结果（相差一个无关紧要的%
\trDeformation）就是 $D_{H,\lambda}$。所要的结论现由\eqref{eq:delbar-gluing}
得出。
\end{proof}
